\section{Behavioral-Coverage Diagnostic}
\label{app:analysiscoverage}

\phantomsection
\textbf{Source-grounded construction and refinement.}\quad
\label{app:analysisconstruction}
Equations~\ref{eq:mechanismledger}--\ref{eq:residualmap} describe the
coverage-checked variant evaluated here, not a retroactive change to the
training-study implementation. Same-task repeats are analyzed together;
oversized packs are split into complete episodes. The ledger retains
source-backed statements and their task contexts independently of successive
drafts. Required behavior remains distinct from observed mastery, and recovered
errors from unresolved failures. A source ID verifies ancestry, not semantic
retention; unavailable analyses remain missing evidence.

Hierarchical merging and source-scoped review form a report whose shared
premises are then factored under a word target. Checks compare its claims with
their cited mechanisms and public contexts. Local edits repair explicit
conflicts; residual completion consolidates only unexpressed relations,
retaining relevant actor, authorization, scope, ordering, and recovery
conditions. Equivalent content is matched to a literal existing passage
rather than appended again. Incidental names and repeated narratives may be
removed, but operative distinctions are retained as separate branches.
The final variant reuses its own compact draft, local repairs, and
internal missing-source findings; it does not receive evaluation references,
judge scores, or the native report's prose. Internal checks remain fallible.

\textbf{Task-linked measurement in the maintained framework.}\quad
Separately from report construction, the data researcher maintains requirement
definitions and sets $J_c$ of real public-eval-task IDs. The controller computes
their performance from recorded outcomes:
\begin{equation}
s_t(c)=\frac{1}{|J_c|}\sum_{i\in J_c}
       \frac{1}{K}\sum_{k=1}^{K}y_{t,i,k},
\label{eq:capability}
\end{equation}
where $K$ is the number of evaluation repeats and $y_{t,i,k}$ the recorded
outcome. Changed definitions or memberships require recomputing baseline and
current results on the same membership. Groups may overlap; composite-task
success does not isolate individual skills. This measurement is not inferred
from summary frequency or used as the diagnostic's semantic-coverage score.
Source-to-case queries expose complete permitted interactions and
cross-checkpoint histories; internal merge ancestry is retained without
promising arbitrary intermediate-node browsing.

\textbf{Inputs and comparison.}\quad
We freeze complete initial public-evaluation corpora: BFCL has 400 indexed tasks
and 1,200 trajectories, $\tau^2$ has 139 tasks and 417 trajectories, and
ACEBench-Agent has 25 tasks and 75 trajectories. Each task has three evaluation
repeats; these are not independent analysis seeds. These Qwen3-8B reasoning
trajectories are separate from the Qwen3-4B training curves. Public/private
denotes researcher access, not the benchmark's official release status.
The native analyst uses \texttt{gpt-6-astra}/high in a fresh session with a
minimal task-and-schema instruction, a 60-minute ceiling, and ordinary coding
and search tools. It can choose its own strategy and has no imposed report word
limit. Its three saved reports are fixed for all candidate comparisons.

An independent AppWorld extension applies the same final protocol to 81 public
tasks with three repeats (243 trajectories) from Qwen3-4B-Instruct-2507.
Task and failure coverage are 89.70\%/66.32\% for the native report and
99.49\%/85.47\% for the refined map, aggregated over 80 and 79 eligible tasks,
respectively. The reports contain 5,412 and 4,522 words.
The three-benchmark historical comparisons below remain unchanged.

The refined-map generator uses \texttt{gpt-5.5}/high, complete public task
evidence and its own intermediate records. It receives neither reference labels
nor the agent report's prose; only the latter's word count sets a soft 85\%
length target. The final reports are 3,637, 3,775, and 2,596 words; $\tau^2$
is shorter than the agent but does not meet that 85\% target. Same-task repeats
are grouped, oversized groups are split into complete episodes, and shared
definitions are deduplicated without truncating observations. The initial ACE
scoped-review pilot supplies system rules and the first user request as extra
context; later BFCL/$\tau^2$ scoped reviews supply all public requests. Leaf
extraction and review receive complete episodes in all three corpora.

\textbf{Fixed automatic evaluation.}\quad
References are independently generated from public trajectories with
GPT-5.5/high, before seeing the candidate reports, and retain source quotations.
The same model grades anonymized report pairs as full, partial, missing, or
contradicted. Full requires preserving the mechanism and its consequential
scope, condition, actor, or order; mere category overlap is insufficient.
Only full contributes to the strict score. First average reference units within
each indexed task, then average tasks; the supplemental metric additionally
averages fixed subset--family groups equally. Repeat-level evidence does not
increase the independent task count. The reference inputs, grading prompt,
batch size (24), model settings, blind seed, and aggregation are unchanged
across candidates. Literal supporting quotations are checked mechanically.

\begin{table}[ht]
\centering\small
\setlength{\tabcolsep}{3.5pt}
\begin{tabular}{lrrrr}
\toprule
Corpus & Agent original & Agent paired & Refined map & Words (agent / map) \\
\midrule
BFCL & 71.29 / 47.65 & 61.49 / 39.40 & 88.08 / 71.31 & 5,134 / 3,637 \\
$\tau^2$ & 86.24 / 69.92 & 78.66 / 61.88 & 97.50 / 80.48 & 4,283 / 3,775 \\
ACE-Agent & 88.38 / 76.57 & 82.46 / 67.46 & 98.24 / 88.36 & 3,337 / 2,596 \\
\bottomrule
\end{tabular}

\caption{Task / failure coverage (\%) and report lengths. ``Original'' and
``paired'' grade the same native report in different judging contexts, not
different native generations.}
\label{tab:analysiscoverage}
\end{table}

\begin{table}[ht]
\centering\small
\begin{tabular}{lrrr}
\toprule
Corpus & Agent original & Agent paired & Refined map \\
\midrule
BFCL & 67.68 / 49.77 & 57.81 / 41.23 & 86.00 / 69.41 \\
$\tau^2$ & 70.95 / 55.27 & 58.17 / 49.83 & 95.05 / 74.02 \\
ACE-Agent & 90.73 / 80.11 & 86.02 / 74.34 & 98.12 / 90.53 \\
\bottomrule
\end{tabular}

\caption{Subset--family-balanced task / failure coverage (\%). Families are a
fixed diagnostic taxonomy, not an exhaustive ontology of benchmark failures.}
\end{table}

\textbf{Evaluation provenance and missingness.}\quad
Sections~\ref{sec:coveragegap} and~\ref{sec:analysiscoverage} use the same final
paired-evaluation results. Earlier scores are retained as historical results,
not pooled with that evaluation or interpreted as a same-input repeatability
test. Saved configurations retain the judge model, prompt, reasoning level,
batch size, and blind seed; the paired comparison report changes, as do some
eligible scoring units. The score difference therefore does not isolate judge
instability, and a downward shift in the baseline is not a method contribution.
Both current and original agent scores are below the refined map's four coverage
metrics in every corpus. Historical grades have slightly different questionable
or unscored units and are not paired causal effect estimates. Within each
comparison the methods share the same eligible units. BFCL scores 8,327 units
and ACE 538; $\tau^2$ retains 3,191 of 3,200. Respectively, 75, 14, and 2 units
are marked questionable and excluded jointly. Original reference annotations
are unavailable for 7, 13, and 1 trajectories. One bounded format-retry pass
and the pre-existing valid-paired-unit recovery rule are used; an interrupted
$\tau^2$ retry is concluded from saved answers without another scoring pass.

Astra/high provides a separate, disagreement-enriched spot audit. BFCL, $\tau^2$,
and ACE have 22/24, 23/24, and 22/24 evaluable items: 21/23/22 references are
supported and one BFCL reference is uncertain. Exact first/second-rater label
agreement for agent/map is 15/22 versus 18/22, 20/23 versus 20/23, and 15/22
versus 20/22. These samples are not global accuracy estimates or human gold.

\textbf{Illustration and counterexamples.}\quad
In the ACE recovery example discussed in Section~\ref{sec:analysiscoverage}, a successful purchase is followed by a capacity-blocked
notification. The agent retrieves messages but reaches its turn limit without
obtaining a deletion target, deleting, or retrying. The native report mentions
inbox-cleanup truncation; the refined map preserves the recovery dependency and
distinguishes interactive authorization from preset deletion. Both graders
support partial versus full coverage. This is not evidence of deliberate
abandonment or of successful recovery.

The method still loses information: it compresses a BFCL card-number-as-card-ID
error and its recovery into a generic identifier warning, and an ACE repeated
date-argument error after TypeError feedback into a parameter-error category.
For BFCL prefix-only deletion of a 31-file list, the second grader downgrades
the map's ``Drafts partial deletion'' from full to partial. A separate purposeful
four-item BFCL case audit retains three evaluable results and one format
failure; it does not modify primary scores or enter the main audit counts.

\textbf{Compute and interpretation.}\quad
Generation costs include reused source extraction and parent stages, not just
the final compression call. The map uses substantially more computation; input
includes cached tokens and output includes reasoning tokens. Unknown usage is
not zero, and token totals across different models are not monetary costs.
Reference construction, scoring and audit requests are evaluation overhead
preserved separately in the study logs.

\begin{table}[ht]
\centering\small
\setlength{\tabcolsep}{4pt}
\begin{tabular}{lrrrr}
\toprule
Corpus / method & Input (M) & Cached (M) & Output (M) & Unknown usage \\
\midrule
BFCL / Agent & 0.876 & 0.766 & 0.017 & 0 \\
BFCL / Map & 43.241 & 0.503 & 8.949 & 20 \\
$\tau^2$ / Agent & 1.327 & 1.213 & 0.016 & 0 \\
$\tau^2$ / Map & 16.192 & 0.593 & 4.077 & 0 \\
ACE-Agent / Agent & 1.462 & 1.363 & 0.012 & 0 \\
ACE-Agent / Map & 1.049 & 0.047 & 0.591 & 0 \\
\bottomrule
\end{tabular}

\caption{Recorded generation tokens in millions; cached input is a subset of
input. Unknown-usage requests are disclosed separately. These totals describe
the studied prototype, not an efficiency-matched comparison.}
\end{table}

All candidate versions and negative results are retained. The original simple
merge did not uniformly outperform the agent; scoped expansion improved detail
but grew long, while the final refinement restored only semantic residuals.
The final candidate does not dominate every intermediate version on every
metric. These corpora were used for method development, not an unseen
confirmation study. The diagnostic supports reduced information loss in this
setting, not semantic losslessness, a same-model/same-budget causal attribution,
the independent value of source lookup, or downstream RL improvement. The
refinement is not retroactively attributed to the historical training runs.

\textbf{Provenance of the motivating cases.}\quad
The coverage scores in Section~\ref{sec:coveragegap} use the final native
strict-coverage grades from the same frozen results as Figure~\ref{fig:analysiscoverage}a.
The archived BFCL case audit uses public trajectory
\texttt{efefe1b89e2fb4e8382f9a5f0c230912}, preserved locally with the diagnostic
sources. Six successful file-deletion calls target the first six entries of
the 31-file listing; the other 25 entries are not targeted. Two subsequent
directory-removal calls report a missing path, yet the final reply claims
complete cleanup. The audit identifies an omitted execution pattern, not an independently
verified final filesystem state or a fully repaired summary.
The archived native excerpt is a generic filesystem finding, not a diagnosis
of this specific trajectory. The complete-report case audit judges the
prefix-only mechanism missing under both raters; the excerpt alone does not
establish absence. Its verbatim text and report fingerprint are retained in
\texttt{bfcl\_prefix\_report\_excerpt.json} beside the local case and audit.
The separate ACE R9 record audit uses the frozen 4B results: 48 planned tasks, 38 trained
tasks, 304 episodes, five updates, zero recorded episode aborts, and 45.56\%
public-eval score. The purchase-workflow revision's allocation is 10 planned,
zero accepted, and unavailable, with no duplicates or batch-alignment drops.
The case is a plan--exposure record audit, not an unlinked-feedback control;
it neither attributes a score change to missing data nor rules out exposure
to earlier versions. The aggregate scores and original cases are summarized in
\texttt{data/evidence\_gaps\_20260918.json}; no new training or judging was
performed for these illustrations.
